\documentclass[a4paper,11pt]{article}

% --- Encodage et langue ---
\usepackage[utf8]{inputenc}
\usepackage[T1]{fontenc}
\usepackage[french]{babel}

% --- Mise en page ---
\usepackage[margin=2.5cm]{geometry}
\usepackage{parskip}
\usepackage{setspace}
\onehalfspacing

% --- Polices ---
\usepackage{lmodern}

% --- Tableaux ---
\usepackage{booktabs}
\usepackage{array}
\usepackage{longtable}

% --- Mathématiques ---
\usepackage{amsmath}
\usepackage{amssymb}

% --- Listes ---
\usepackage{enumitem}

% --- Liens ---
\usepackage[colorlinks=true, linkcolor=blue!60!black, urlcolor=blue!60!black]{hyperref}

% --- En-tête / pied de page ---
\usepackage{fancyhdr}
\pagestyle{fancy}
\fancyhf{}
\fancyhead[L]{\small Cahier des charges --- Mesure de la marche V2026}
\fancyhead[R]{\small Robin \& Thoma}
\fancyfoot[C]{\thepage}
\renewcommand{\headrulewidth}{0.4pt}

% --- Code inline ---
\usepackage{xcolor}
\definecolor{codebg}{gray}{0.95}
\newcommand{\code}[1]{\colorbox{codebg}{\texttt{#1}}}

% ======================================================================
\begin{document}

% --- Page de titre ---
\begin{titlepage}
    \centering
    \vspace*{3cm}
    {\Huge\bfseries Cahier des charges\par}
    \vspace{0.8cm}
    {\LARGE Mesure de la marche V2026\par}
    \vspace{1.5cm}
    {\Large Robin \& Thoma\par}
    \vspace{2cm}
    {\large Projet interdisciplinaire HES-SO Master\\
    Printemps 2026\par}
    \vfill
\end{titlepage}

\tableofcontents
\newpage

% ======================================================================
\section{Introduction}

Ce document décrit le projet interdisciplinaire «~Contrôle Horloger Intelligent~» dans sa version 2026. Il s'inscrit dans la continuité des travaux menés en 2025 et vise à aboutir à un système capable de mesurer automatiquement la marche d'une montre mécanique par voie optique.

Le projet n'est pas terminé. Ce cahier des charges pose le cadre de ce qui a été fait, de ce qui reste à faire, et des objectifs que nous cherchons à atteindre d'ici la fin du semestre.

% ======================================================================
\section{Contexte et problématique}

\subsection{Qu'est-ce que la marche~?}

La \textbf{marche} d'une montre mécanique est l'écart entre sa fréquence réelle de battement et sa fréquence nominale, exprimé en secondes par jour. Concrètement, une montre qui «~avance~» de 5~secondes par jour a une marche de $+5$~s/j. Les normes COSC imposent une tolérance de \textbf{$-4$ à $+6$~s/j} pour certifier un chronomètre.

\subsection{Comment la mesure-t-on aujourd'hui~?}

La méthode standard dans l'industrie est \textbf{acoustique}~: un microphone capte les «~tic-tac~» du mécanisme et un appareil (Witschi, Timegrapher) en déduit la fréquence. C'est fiable, rapide, bien maîtrisé.

\subsection{Pourquoi une approche optique~?}

L'idée de ce projet est d'explorer une alternative \textbf{par imagerie stroboscopique}. Le principe~: on éclaire le balancier d'une montre avec un flash à une fréquence très proche de sa fréquence propre. En résulte un effet de «~ralenti apparent~» --- comme une roue qui semble tourner lentement sous un éclairage néon. En analysant ce mouvement apparent image par image, on peut remonter à la fréquence réelle du balancier et donc calculer la marche.

Cette approche présente un intérêt pour les cas où la méthode acoustique est inadaptée (mouvements silencieux, environnement bruité) et ouvre la porte à des analyses visuelles complémentaires (détection d'anomalies, état des composants).

% ======================================================================
\section{Ce qui a été réalisé en 2025}

L'édition 2025 du projet a posé les fondations matérielles et logicielles du système. Voici ce dont nous héritons.

\subsection{Station d'acquisition}

Une station physique a été conçue et fabriquée (pièces imprimées 3D, fichiers CAO en STEP disponibles). Elle comprend~:

\begin{itemize}
    \item Un \textbf{dôme} d'éclairage avec LED commandée par un microcontrôleur RPi Pico
    \item Un \textbf{support montre} ajustable
    \item Des \textbf{bras de réglage X/Y} pour le positionnement de la caméra
    \item Une \textbf{caméra GigE Vision} SVS EXO273CGE ($1440 \times 1080$, monochrome)
\end{itemize}

Le flasher est piloté par liaison série (USB-UART, 115\,200~baud). Le protocole est simple~: des commandes texte au format \code{commande:valeur;} permettent de régler la période du flash, la durée de l'éclair, l'exposition, etc.

\subsection{Interface de contrôle (GUI)}

L'équipe 2025 a développé une application graphique complète sous PyQt6 pour piloter manuellement le système~:

\begin{itemize}
    \item Contrôle en temps réel des 6~paramètres du flasher via sliders et champs texte
    \item Connexion caméra avec prévisualisation live et réticule configurable
    \item Mode burst (rafale de 2 à 60~images) pour capturer des séquences du balancier
    \item Sauvegarde/chargement de configurations (format \code{.ini})
    \item Logs série en temps réel
    \item Mode test (port série fictif + caméra factice) pour développer sans matériel
\end{itemize}

\subsection{Classification des images}

Un classifieur par analyse de signal a été développé. Il prend 4~images consécutives en entrée et identifie la situation de la dernière image en 4~classes~:

\begin{table}[h]
    \centering
    \begin{tabular}{c l l}
        \toprule
        \textbf{Classe} & \textbf{Ce qu'on observe} & \textbf{Ce que ça signifie} \\
        \midrule
        0 & 1 pic d'intensité & Le balancier est quasi-statique dans le champ \\
        1 & 2 pics (faible $\to$ fort) & Il s'approche de sa position d'équilibre \\
        2 & 2 pics (fort $\to$ faible) & Il s'éloigne de sa position d'équilibre \\
        3 & 3 pics & Il traverse l'équilibre --- c'est le «~saut~» \\
        \bottomrule
    \end{tabular}
    \caption{Classes de classification du signal du balancier.}
    \label{tab:classes}
\end{table}

La méthode repose sur une soustraction de la médiane des images, l'extraction d'un profil d'intensité perpendiculaire au signal, et une détection de pics (\code{scipy.signal.find\_peaks}).

\subsection{En résumé~: ce qui marchait, ce qui manquait}

À la fin de 2025, on disposait d'un système \textbf{interactif}~: un opérateur pouvait capturer des images, les classifier, et observer visuellement si le flash était synchronisé. Mais tout était \textbf{manuel}. Il n'y avait pas de boucle de rétroaction automatique, pas de calcul de marche, pas de validation métrologique.

% ======================================================================
\section{Objectifs de la version 2026}

Notre ambition pour 2026 est de transformer ce prototype interactif en un \textbf{outil de mesure autonome}. Concrètement~:

\begin{enumerate}
    \item \textbf{Automatiser la synchronisation} --- Le système doit pouvoir ajuster lui-même la fréquence du flash jusqu'à se caler sur le balancier, sans intervention humaine.

    \item \textbf{Calculer la marche} --- Une fois synchronisé, le système doit mesurer la fréquence apparente du balancier et en déduire la marche en s/j.

    \item \textbf{Valider la mesure} --- Il faut pouvoir comparer nos résultats avec une référence (si possible) et caractériser la précision du système (répétabilité, justesse, incertitude).

    \item \textbf{Bonus} --- Optimisation de la classification~? GUI~?
\end{enumerate}

% ======================================================================
\section{Description fonctionnelle}

\subsection{Pilotage du flasher}

On reprend le protocole série existant (RPi Pico). Les paramètres clés sont~:

\begin{table}[h]
    \centering
    \begin{tabular}{l l l}
        \toprule
        \textbf{Paramètre} & \textbf{Rôle} & \textbf{Plage} \\
        \midrule
        \code{trig\_off}   & Période du flash (détermine la fréquence) & $1\,000$ -- $1\,000\,000$~µs \\
        \code{flash\_on}   & Durée de l'éclair               & $1$ -- $100\,000$~µs \\
        \code{flash\_off}  & Temps entre fin d'éclair et prochain trigger & $1$ -- $100\,000$~µs \\
        \code{trig\_expo}  & Temps minimal d'exposition       & $0$ -- $200$~µs \\
        \code{trig\_shift} & Décalage temporel du trigger     & $0$ -- $100\,000$~µs \\
        \bottomrule
    \end{tabular}
    \caption{Paramètres de commande du flasher.}
    \label{tab:flasher-params}
\end{table}

Le \code{trig\_off} est le paramètre central~: c'est lui qu'on ajuste pendant la synchronisation. Sa valeur initiale dépend du calibre de la montre~:

\begin{table}[h]
    \centering
    \begin{tabular}{r r r}
        \toprule
        \textbf{Calibre (A/h)} & \textbf{Fréquence nominale (Hz)} & \textbf{\code{trig\_off} initial (µs)} \\
        \midrule
        18\,000 & 2.5 & 400\,000 \\
        21\,600 & 3.0 & 333\,333 \\
        25\,200 & 3.5 & 285\,714 \\
        28\,800 & 4.0 & 250\,000 \\
        36\,000 & 5.0 & 200\,000 \\
        \bottomrule
    \end{tabular}
    \caption{Valeurs initiales de \code{trig\_off} selon le calibre.}
    \label{tab:calibres}
\end{table}

\subsection{Acquisition d'images}

Points importants~:

\begin{itemize}
    \item Résolution~: à déterminer.
    \item Les images sont très sombres. On applique une \textbf{normalisation CLAHE} (\emph{Contrast Limited Adaptive Histogram Equalization}) pour rendre les images exploitables pour l'affichage. Son utilisation avant la classification est en cours de test.
\end{itemize}

\subsection{Classification}

On reprend le classifieur de 2025 avec des ajustements. Le principe reste le même~: à partir de 4~images, on soustrait la médiane, on extrait un profil d'intensité perpendiculaire, et on compte les pics pour déterminer la classe (0, 1, 2, 3 ou $-1$ si inclassable).

\subsection{Synchronisation automatique}

C'est la pièce maîtresse de la version 2026. L'idée~: ajuster automatiquement \code{trig\_off} pour que la fréquence du flash colle à celle du balancier. Quand c'est le cas, on observe un cycle stable de classes ($0 \to 1 \to 3 \to 2 \to 0 \to \ldots$).

\medskip
\noindent\textit{(Phase de test)}

\medskip
L'algorithme envisagé fonctionne en deux phases~:

\begin{enumerate}
    \item \textbf{SEARCH} --- On balaie \code{trig\_off} par pas grossiers (de l'ordre de 100~µs) autour de la valeur nominale, dans les deux directions. Le but est de trouver la zone où des classes 1, 2 ou 3 apparaissent. Sans cette phase de recherche large, on peut facilement passer à côté si la fréquence réelle de la montre est un peu éloignée de la nominale.

    \item \textbf{LOCK} --- Une fois la zone trouvée, on passe en pas fins (de l'ordre de 10~µs) pour stabiliser le taux de classes~3 dans une fenêtre glissante. Le verrouillage est confirmé après plusieurs itérations stables.
\end{enumerate}

Ce mécanisme est encore en cours de mise au point. Les premiers essais ont révélé des problèmes de convergence (balayage unidirectionnel qui dérive sans jamais trouver, faux verrouillage sans vraie classe~3). L'approche bidirectionnelle et les critères de verrouillage doivent encore être validés sur le matériel.

\subsection{Calcul de la marche}

Une fois le flash synchronisé, on acquiert des images pendant une durée configurable (10~secondes par défaut) et on mesure la fréquence apparente du balancier.

Le principe mathématique~:

\begin{equation}
    f_{\text{réelle}} = f_{\text{flash}} + f_{\text{apparente}} \times \text{signe}
\end{equation}

\begin{equation}
    \text{marche} = \frac{f_{\text{réelle}} - f_{\text{nominale}}}{f_{\text{nominale}}} \times 86\,400 \;\text{s/j}
\end{equation}

La fréquence apparente se déduit du nombre de cycles complets observés (transitions $0 \to 1 \to 3 \to 2 \to 0$) sur la durée de mesure. Le signe (avance ou retard) est déterminé par l'ordre d'apparition des classes~1 et~2.

% ======================================================================
\section{Contraintes techniques}

\subsection{Matériel existant}

\begin{table}[h]
    \centering
    \begin{tabular}{l l l}
        \toprule
        \textbf{Composant} & \textbf{Modèle / Détail} & \textbf{Remarque} \\
        \midrule
        Caméra    & SVS EXO273CGE, $1440 \times 1080$, mono & GigE Vision, compatible Aravis \\
        Éclairage & LED + RPi Pico              & Firmware série existant \\
        Station   & Pièces imprimées 3D          & CAO livrée (STEP) \\
        Liaison   & USB-UART                     & Protocole texte \code{cmd:val;} \\
        \bottomrule
    \end{tabular}
    \caption{Matériel du banc de mesure.}
    \label{tab:materiel}
\end{table}

\subsection{Logiciel}

\begin{table}[h]
    \centering
    \begin{tabular}{l l}
        \toprule
        \textbf{Critère} & \textbf{Choix} \\
        \midrule
        Langage              & Python $\geq$ 3.10 \\
        Pilote caméra        & Aravis via PyGObject (GI v0.8) \\
        Traitement d'image   & OpenCV, NumPy \\
        Analyse de signal    & SciPy (\code{find\_peaks}) \\
        Communication série  & pyserial \\
        OS cible             & macOS (ARM/Intel), Linux \\
        \bottomrule
    \end{tabular}
    \caption{Choix logiciels.}
    \label{tab:logiciel}
\end{table}

\vfill
\begin{center}
    \rule{0.5\textwidth}{0.4pt}\\[0.5em]
    \textit{Projet interdisciplinaire HES-SO Master --- Contrôle Horloger Intelligent, printemps 2026.}
\end{center}

\end{document}
